\section{Checking the field and spectrum endpoints}\label{app:tech}
The finite cell model uses the equal-area disk self term \cite{richmond}. For square cells of side $h$, set $z=kh/\sqrt\pi$. The Green entries away from the diagonal are $\mathrm i\pi zJ_1(z) H_0^{(1)}(k d)/2$ at centre-to-centre distance $d$. The diagonal is $\mathrm i\pi z H_1^{(1)}(z)/2-1$. A port outside the region uses the off-diagonal rule. We retain the incident field, complex datum normalization and Green representation supplied with each array. In the main large-domain comparison the bound applies to the exact represented padded Fourier symbol and the original receiver and illumination coefficients. Known exterior pixels can be eliminated through the corresponding block system. For example, with known contrast $M_h$ and perturbations only on the selected cells $T_c$, the active-field equation is obtained by applying $T_c^*(I-GM_h)^{-1}$ to $B+G T_cq_c$. This accounts for every interaction with the fixed host. No coarse eigenvectors are taken as a complete fine-grid basis.

An interval matrix is evaluated by midpoint--radius arithmetic with outward rounding \cite{arb}. To count a whole nonsingular state spectrum we use an approximate frame $V$ as a proposal, verify $\lambda_{\min}(V^*V)>0$, and enclose $V^*KV$ by signed, separated intervals or a one-sided complement estimate. For the direct test at $t$ it is enough to exhibit $n+j$ independent columns with positive restricted form. On the prospective small boxes a binary64 calculation bounds complex dot-product, assembly, frame and residual errors before an online decision. Each positive result was also replayed directly by evaluating Gram products of the proposed columns in ball arithmetic at the original endpoints. This second computation avoids assumptions about the proposing eigensolver. Candidate forward solves are upper endpoints; a small observed floating residual alone cannot justify an exclusion or inclusion. In the signed $32^2$ experiment a full-state residual and a verified inverse-defect bound enclose the forward response. No positivity of the material was used.

A rational singular example checks the counting convention. In excess-current coordinates take one active cell, $h=W=1,V=0,T=5/4$, two sources, $E_c=(3/4,0)$, and two receivers with $P=(3/4,0)^\top$, $y=\diag(-5/4,2)$. Then $K=L=0$, $C=\diag(1,4)$. The count integral is five and the admissible contrast perturbation $-1$ attains squared residual five. Here $\mathcal A-tP_s=\diag(0,1-t,4-t)$. The uncompressed strict count $\np\ge1+j$ can detect only $j=1$ and $t<1$. Removing the exact common null state restores the full two-source spectrum in the integrated count. This is an algebraic control for a persistent zero. The large radiating examples use a resolved numerical spectrum.

For the full-domain branch calculation write the baseline quadratic at $[0,b]$ in units of its positive sector multiplier. The uncompressed state is $I-b\Herm G+S^*S/\mu$. A finite subset of a Fourier convolution is a compression of its padded embedding, so this operator is bounded below by $(1-b\max\Re\widehat G)I+S^*S/\mu$. When the scalar floor is positive, its inverse is computable by a receiver-size system. We keep the exact simultaneous linear and constant terms of the original base quadratic. After replacing the state block by this lower block, local changes to the endpoint rows have exactly the form \eqref{eq:update}. Arb Fourier inversion encloses the required submatrix of $G$, the two row Grams, all sources and receiver products, and the state and source eigenvalue bounds. For each patch the reduced inertia is checked on all $2k_p$ coordinates, and the positive tail is rounded down using a near-unitary frame and its verified defect. Unknown background pixels have entered the calculation before this reduction. A positive base floor is an explicit sufficient condition for this implementation.

The conventional perturbation comparison uses the same background, active coordinates and endpoints. In the disk example let $T$ denote the exact active Green operator about the candidate material. Interval linear solves bound the two-norm from below and above. An upper perturbation $wI$ respecting the box gives $n\rho(wT)^4\ge|\tr((wT)^4)|$ by the complex eigenvalue trace identity; it does not require normality. When floating construction of tabulated material endpoints rounds inward, the proof for an ideal real interval uses its largest guaranteed contained width. This accounts for less than $4.5\,10^{-16}$ at the quoted disk width.

In the structured 49-cell problem, put $U=[S^\top,B]$. We verified a left inverse $LU=I$ for the corresponding joint field frame in both the structured-pool support comparison and the separate applicable flow geometry. For any complex matrix $Y$, $T_0=L^\top\bigl(\begin{smallmatrix}0&Y\\Y^\top&0\end{smallmatrix}\bigr)L$ is symmetric and has $ST_0B=Y$. Thus support and unbounded reciprocity alone admit every measured matrix on these arrays. Such a $T_0$ need not be produced by a passive diagonal material. It is the missing constitutive information that the sector supplies.

A six-cell example checks how a sound branch test handles actual ambiguity. We reused two known numerical materials and an existing three-illumination, single-receiver measurement, without repeating the root search. Independent 128-bit evaluations of the original full-wave equations bound their residuals by $1.603\,10^{-15}$ and $1.415\,10^{-15}$, compared with $r=2^{-30}\|Y\|_F>7.6340\,10^{-11}$. Their RMS material separation exceeds 1.15879. Disjoint real boxes of half-width 0.05 around their respective centres must both survive. This proves compatibility of distinct materials with the declared observation ball. It does not assign algebraic multiplicity to the rounded central reading.

\section{Scope of the blind-box comparison}\label{app:limits}
For nonsingular $K$ diagonalize its negative and positive parts and write a physical displacement as $J-J_0=Nc+Pv$ in energy-normalized bases, with $J_0=K^{-1}L$. Then $D\succeq\mathcal H+v^*v-c^*c$. Individual material rows require
\[
 [(I-XG)P]_i v=[X(B+GJ_0)-J_0]_i-[(I-XG)N]_i c .
\]
The nonzero source-dependent term must be retained. If these rows imply $v^*v+M_Y\succeq c^*\Pi c$ on the whole tested box for an orthogonal projection of rank $c_0$, only $\nu-c_0$ unconstrained coordinates remain. Interlacing gives $\Tail_{\nu-c_0}(\mathcal H-M_Y)$ as a lower bound. For the fixed 13-mode state form on the six $96^2$ arrays, a verified upper estimate on the attainable gain of separately weighted rows rules out every selection of at most 59 cells as a route to a cap-one exclusion, including the tested signed centre charges. This is a limit on that local-row construction. It neither constructs a low-peak material nor bounds methods that exploit a different material hypothesis. Equations \eqref{eq:row} and \eqref{eq:border} change all of the relevant rows when a new spatial interval is specified, which is why the two results are consistent.

In the large-domain study, contrast and total incident fluence are held fixed when varying the number of sources. Measurements can add angular information until the port fields repeat combinations within tolerance. On a half aperture the restrictions of passivity and optical power can constrain receiver-invisible currents that still radiate. On a dense full ring the corresponding remaining modes can radiate extremely weakly. Their finite number across sources, rather than a large positive optical eigenvalue, can still strengthen a distance estimate. Neither mechanism makes spatial recovery unique. At a single voxel $G={\rm i}$, $S=B=1$, the midpoint of $J(0)$ and $J(1)$ would require complex contrast $2/5+{\rm i}/5$. Even a scalar real material interval can have a nonconvex set of full-wave responses.
